\section{常见故障与治理}

\subsection{故障共性}

\begin{warningbox}{训练中的泛化失败}
即使在训练环境中获得高 reward，agent 仍可能在稍有改动的初始状态下崩溃。原因通常是 replay buffer 过度依赖一小批高 reward 转移，难以覆盖状态空间。可以借助 entropy bonus 或多任务 replay 来缓解。
\end{warningbox}

\begin{knowledgebox}{治理关注点}
\begin{itemize}
    \item Policy drift：需保留 dev 策略样本作为窗内基准。
    \item Reward hacking：强化学习中 reward shaping 要谨慎，不可通过简单修改 reward 让 agent 找到 shortcut。
    \item 数据透明：记录每次更新使用的数据 slice，并定期与业务方核对行为日志。
\end{itemize}
\end{knowledgebox}

\subsection{干预机制}

\begin{importantbox}{Safe Interruptibility 与人为干预}
在训练过程中，保持 interruptibility 意味着当人类检测到异常时，可以安全终止或重置 agent，而不会引发奖励爆炸。实践中会设置 \texttt{kill switch}，同时记录最后一次 TD error 以帮助定位原因。
\end{importantbox}

\begin{warningbox}{过度干预带来的训练碎片化}
频繁中断训练会导致 replay buffer 中数据来自不同策略，破坏 off-policy 假设。建议在干预后执行 warm-start offline training，确保新策略有充分时间收敛。
\end{warningbox}

\subsection{本章小结}

治理层面需平衡自动训练与人为监控，干预机制应以 guardrail 为核心，避免重复训练碎片化。

